% =====================================================================
%  ALANYA — Dossier de conception
%  Exportation des médias, sauvegarde et restauration de compte
%
%  Compilation :  latexmk -pdf sauvegarde-exportation-medias.tex
% =====================================================================
\documentclass[11pt,a4paper,oneside]{report}

\usepackage{alanya-conception}
\setAlanyaDocTitre{Sauvegarde et exportation}

\hypersetup{
  colorlinks=true,
  linkcolor=AlanyaIndigoStrong,
  urlcolor=AlanyaIndigoStrong,
  pdftitle={Alanya - Exportation des medias, sauvegarde et restauration},
  pdfauthor={Equipe Alanya},
  pdfsubject={Conception - export local, sauvegarde Google Drive, migration d'appareil}
}

\begin{document}

% =====================================================================
\begin{titlepage}
\thispagestyle{empty}
\begin{tikzpicture}[remember picture, overlay]
  \fill[AlanyaIndigo]
    (current page.north west) rectangle ([yshift=-6.4cm]current page.north east);
  \fill[AlanyaIndigoDark]
    ([yshift=-6.4cm]current page.north west) rectangle ([yshift=-6.0cm]current page.north east);
\end{tikzpicture}

\vspace*{1.1cm}
{\color{white}\sffamily
\begin{center}
  {\small\bfseries DOSSIER DE CONCEPTION}\\[6pt]
  {\footnotesize\color{AlanyaContainer} Document de travail --- à valider avant développement}
\end{center}}

\vspace{3.4cm}

\begin{center}\sffamily
  {\color{AlanyaLine}\rule{0.6\linewidth}{0.6pt}}\\[20pt]
  {\fontsize{30}{34}\selectfont\bfseries\color{AlanyaIndigoDark} Exportation,}\\[4pt]
  {\fontsize{30}{34}\selectfont\bfseries\color{AlanyaIndigoDark} sauvegarde}\\[4pt]
  {\fontsize{30}{34}\selectfont\bfseries\color{AlanyaIndigoDark} et restauration}\\[10pt]
  {\large\color{AlanyaGray} Ce qui est déjà sur le téléphone n'a pas à repasser par le réseau}\\[6pt]
  {\normalsize\itshape\color{AlanyaGray}
   Archive locale, sauvegarde chiffrée sur Drive, migration d'appareil}\\[22pt]
  {\color{AlanyaLine}\rule{0.6\linewidth}{0.6pt}}
\end{center}

\vspace{1.0cm}

\begin{center}
  \badgeMarque[3.4]
\end{center}

\vfill

\begin{center}\sffamily
  \begin{tabular}{@{}r@{\quad}l@{}}
    \thd{Périmètre}     & Talky (Flutter, Android) \textbullet{} Alanya-Backend\\[3pt]
    \thd{Statut}        & Conception --- aucune ligne de code écrite\\[3pt]
    \thd{Point d'entrée} & Profil \textrightarrow{} Mes médias\\[3pt]
    \thd{Dépend de}     & Partitions de médias (rétention serveur, \code{410 MEDIA\_EXPIRED})\\[3pt]
    \thd{Plateforme}    & Android d'abord --- iOS et iCloud hors périmètre\\[3pt]
    \thd{Date}          & \today\\
  \end{tabular}
\end{center}

\vspace{0.8cm}
\begin{center}\color{AlanyaGray}\small\sffamily
  Document confidentiel --- diffusion restreinte
\end{center}
\end{titlepage}

% =====================================================================
\chapter*{Avant-propos}
\addcontentsline{toc}{chapter}{Avant-propos}
\markboth{Avant-propos}{}
\thispagestyle{plain}

Ce dossier répond à une note manuscrite posée en marge de la maquette de l'écran
\emph{Mes médias}~:

\begin{quote}\itshape
Nous devons travailler à s'assurer que les données sont bien stockées sur le téléphone des
inscrits à la plateforme. Cette exportation doit exploiter cet aspect, ceci nous permettra
d'économiser notre bande passante mais aussi les data de l'inscrit. La seule situation où
il y a une réelle exportation, c'est lorsqu'il ouvre son compte sur un nouveau téléphone~;
dans ce cas les données stockées sur le drive de Google sont transférées sur le nouveau
téléphone.
\end{quote}

Deux exigences s'y lisent, et elles ne se recouvrent pas. La première est une règle
d'économie~: une exportation ne doit rien retélécharger de ce que le téléphone possède
déjà. La seconde est un scénario~: la migration vers un nouvel appareil, où un transfert
réel a lieu, depuis Google Drive.

\begin{note}[Comment lire ce document]
La \textbf{partie~I} décrit ce que l'inscrit voit et ce qu'on lui promet. Elle se valide
sans lire la suite. La \textbf{partie~II} décrit comment on le construit, fichier par
fichier. Le chapitre~\ref{ch:decisions} récapitule toutes les décisions arrêtées et sert
de référence rapide.
\end{note}

\tableofcontents

% =====================================================================
\part{Ce que l'inscrit voit}
% =====================================================================

\chapter{Le problème posé}
\label{ch:probleme}

\section{Deux horloges qui ne tournent pas à la même vitesse}

Depuis la mise en place du stockage partitionné, Alanya tient deux règles de conservation
distinctes, et c'est de leur écart que naît tout ce document.

\begin{center}
\begin{tabularx}{\linewidth}{@{}l l X@{}}
\toprule
\thd{Où} & \thd{Durée} & \thd{Pourquoi}\\
\midrule
Sur nos serveurs & 30 jours & Le coût de stockage est le nôtre. Les partitions
journalières font tomber la plus ancienne chaque nuit.\\
Sur le téléphone & sans limite & Le stockage appartient à l'inscrit. Tout média
téléchargé reste jusqu'à désinstallation ou changement de compte.\\
\bottomrule
\end{tabularx}
\end{center}

\begin{constat}
La valeur de rétention n'est pas codée en dur dans l'application. \code{MediaExpiryPolicy}
l'apprend du champ \code{retentionDays} porté par chaque réponse \code{410 Gone}, et la
mémorise. Elle vaut 365~jours pendant la mise en service du stockage partitionné, 30
ensuite. Toute la présente conception s'appuie sur cette valeur apprise, jamais sur une
constante.
\end{constat}

Passé la fenêtre serveur, un média n'existe plus qu'en un seul exemplaire au monde~: celui
posé sur le téléphone de l'inscrit. Il n'y a plus de filet. Une casse, un vol, un
formatage, un changement de téléphone, et la photo est perdue --- définitivement, sans que
personne chez Alanya ne puisse rien y faire.

\section{Ce que la note demande}

L'exigence d'économie n'est pas un détail d'optimisation, c'est une contrainte de
conception. Un inscrit qui exporte les 47~médias de mars 2026 les possède déjà tous sur son
téléphone. Les redemander au serveur reviendrait à payer deux fois~: notre bande passante
sortante, et le forfait data de l'inscrit, pour des octets qui sont déjà à trente
centimètres de l'écran.

\begin{constat}
\code{ChatDao.watchLocalMedia} ne liste que les messages dont \code{localMediaPath} est
renseigné, c'est-à-dire les fichiers réellement présents sur le disque. L'écran
\emph{Mes médias} est donc déjà, par construction, une vue de ce que le téléphone possède
--- et fonctionne hors connexion. L'exportation hérite gratuitement de cette propriété.
\end{constat}

\section{La réponse retenue : deux mécanismes, jamais confondus}

Un même bouton ne peut pas servir deux besoins aussi différents que « emporter mes photos
de mars ailleurs » et « retrouver mon compte sur mon nouveau téléphone ». Les confondre
mène à la pire des situations~: deux mécanismes qui écrivent au même endroit, et plus
personne ne sait lequel détient ses données le jour où il en a besoin.

\begin{center}
\begin{tabularx}{\linewidth}{@{}l X X@{}}
\toprule
 & \thd{Sauvegarde} & \thd{Exportation}\\
\midrule
Déclenchement & automatique, selon la fréquence choisie & l'inscrit appuie sur un bouton\\
Contenu       & la base triée, plus les préférences retenues & les médias d'une période\\
Format        & interne, chiffré, illisible hors d'Alanya & ouvert, en clair, lisible partout\\
Destination   & dossier \code{Alanya} du Drive de l'inscrit & partage système, Téléchargements, ou Drive\\
Sert à        & retrouver son compte sur un nouveau téléphone & sortir ses souvenirs d'Alanya\\
Coût réseau   & quelques Mo & zéro, sauf récupération explicitement acceptée\\
\bottomrule
\end{tabularx}
\end{center}

\begin{alerte}[Conséquence assumée]
La sauvegarde ne contient pas les médias. Sur un nouveau téléphone, l'inscrit retrouve ses
conversations, mais les médias de plus de 30~jours ne reviennent pas~: ils n'étaient ni sur
le serveur, ni dans la sauvegarde. Ce choix évite de déverser plusieurs gigaoctets dans un
quota Google de 15~Go, mais il déplace l'exportation~: elle n'est plus un confort, elle est
\textbf{le seul filet} pour les médias anciens. C'est pourquoi l'archive exportée est
elle-même relisible par la restauration (chapitre~\ref{ch:restauration}).
\end{alerte}

\chapter{Exporter une période}
\label{ch:export}

\section{Le parcours}

L'inscrit arrive sur \emph{Mes médias}, resserre par discussion, par période, par famille
de média, puis appuie sur \textbf{Exporter cette période}. Une feuille s'ouvre et annonce
trois choses avant tout engagement~: combien d'éléments, quel poids, et ce qui manque.

\begin{enumerate}[leftmargin=*]
  \item \thd{Récapitulatif.} « 47~éléments \textbullet{} 128~Mo \textbullet{} mars 2026
        \textbullet{} Ama Nguema ». Le périmètre exporté est exactement celui que la grille
        affiche~: aucune surprise, ni en plus ni en moins.
  \item \thd{Éléments manquants}, s'il y en a (voir §\ref{sec:manquants}).
  \item \thd{Destination.} Partager \textbullet{} Enregistrer dans Téléchargements
        \textbullet{} Déposer sur mon Drive.
  \item \thd{Assemblage} en tâche de fond, progression en notification.
  \item \thd{Remise} de l'archive selon la destination choisie.
\end{enumerate}

\section{Ce que contient l'archive}

Des fichiers ouverts, rangés comme l'inscrit s'y attend, plus deux fichiers de contexte.

\begin{lstlisting}[style=Alanya, language={}]
alanya-export-241030112-2026-03-01_2026-03-31.zip
  index.html                     <- planche-contact, ouvrable sans Alanya
  manifest.json                  <- contexte lisible par machine
  Alanya Images/
    2026-03-14_ama-nguema_0001.jpg
  Alanya Videos/
  Alanya Audio/
  Alanya Documents/
\end{lstlisting}

\begin{constat}
L'arborescence reprend celle que \code{AlanyaMediaExportService} pose déjà dans
\code{Download/Alanya/}~: \code{Alanya Images}, \code{Alanya Videos},
\code{Alanya Documents}. Un inscrit qui a déjà exporté des médias un par un retrouve la
même organisation dans le zip. Seul \code{Alanya Audio} est nouveau.
\end{constat}

Le \code{manifest.json} est ce qui distingue une archive utile d'un tas de fichiers. Un an
plus tard, un nom de fichier ne dit ni qui a envoyé la photo, ni dans quelle conversation,
ni quand. Le manifeste porte tout cela --- et, accessoirement, les identifiants dont la
restauration a besoin pour remettre chaque fichier en face de son message.

\section{Les éléments manquants}
\label{sec:manquants}

Un média peut manquer à l'appel pour deux raisons qui n'ont rien à voir, et l'écran doit
les distinguer~: promettre un téléchargement qui échouera serait pire que de ne rien
promettre.

\begin{center}
\begin{tabularx}{\linewidth}{@{}l X l@{}}
\toprule
\thd{Cause} & \thd{Situation} & \thd{Traitement}\\
\midrule
Jamais téléchargé &
Reçu, mais le téléchargement automatique était coupé, ou il a échoué. Encore présent sur le
serveur. &
Récupérable\\
Purgé du serveur &
Plus vieux que la rétention. Le serveur répond \code{410}. &
Irrécupérable\\
\bottomrule
\end{tabularx}
\end{center}

La feuille d'export annonce donc, le cas échéant~: « \emph{5~éléments manquants. 3~peuvent
être récupérés (40~Mo). 2~ne sont plus disponibles.} » avec une case à cocher pour le
téléchargement.

\begin{note}[Pourquoi cette entorse à la règle d'économie]
Récupérer avant d'archiver consomme la bande passante que la note veut économiser. L'entorse
est acceptée parce qu'elle est \textbf{explicite, chiffrée d'avance et non cochée par
défaut}~: l'inscrit décide, en connaissant le coût. Ce qui est proscrit, c'est le
téléchargement silencieux --- pas le téléchargement choisi.
\end{note}

Les irrécupérables ne disparaissent pas pour autant~: ils sont inscrits dans le manifeste
avec leur motif. Une archive incomplète qui le dit vaut infiniment mieux qu'une archive
incomplète qui se tait~: l'inscrit ne découvrirait le trou que le jour où il chercherait
la photo.

\section{Le dépôt sur Drive}

Déposer l'archive sur le Drive personnel de l'inscrit règle le problème ouvert par la
sauvegarde sans médias~: ses fichiers anciens survivent à son téléphone, sans qu'Alanya
porte le coût d'une sauvegarde de plusieurs gigaoctets. L'archive va dans le dossier
visible \code{Alanya/}, celui-là même qui accueille les sauvegardes, et la restauration
sait la relire.

\begin{alerte}[À dire clairement dans l'écran]
Cette archive est \textbf{en clair}. C'est sa raison d'être~: elle doit s'ouvrir sans
Alanya, sur n'importe quel ordinateur, dans dix ans. Mais cela signifie que ces photos sont
lisibles dans le Drive de l'inscrit, par Google et par toute application à laquelle il a
donné accès à son Drive. L'écran doit le dire avant le dépôt, pas dans les conditions
d'utilisation.
\end{alerte}

\chapter{Sauvegarder et restaurer}
\label{ch:sauvegarde-vue}

\section{Ce qui est sauvegardé}

La base locale \code{talky\_chat.sqlite}, moins ce qui n'a aucun sens sur un autre téléphone,
plus une poignée de préférences qui n'y sont pas. Cette formulation est plus lourde que
« la base et rien d'autre », mais elle est exacte, et la différence porte sur des choses que
l'inscrit remarquerait.

\subsection*{Une objection à traiter d'emblée}

Le serveur \textbf{n'efface jamais les messages}.

\begin{constat}
\code{dataRetentionService.js} purge les statuts expirés (7~jours), l'historique d'appels
(12~mois), le journal de connexions (90~jours), les jobs en échec, les appareils révoqués.
La table \code{message}, elle, n'apparaît nulle part dans ces purges. Seuls les
\emph{fichiers} médias sont purgés, par \code{mediaRetention.js}.
\end{constat}

Autrement dit, sur un téléphone neuf, l'inscrit qui se reconnecte récupérerait de toute façon
ses conversations, ses messages et ses contacts en resynchronisant. Alors pourquoi
sauvegarder ce que le serveur possède déjà~?

\textbf{Parce que la resynchronisation se paie exactement dans la monnaie que la note protège.}
L'application pagine~: \code{getMessages} descend 50~messages à la fois, et
\code{loadOlderMessages} remonte l'historique par tranches de 30. Reconstruire trois ans de
conversations de cette façon, ce sont des milliers d'allers-retours, notre bande passante
sortante et le forfait data de l'inscrit --- pour des octets qu'une archive de quelques
mégaoctets aurait livrés en une seule fois.

\begin{note}[La bonne façon de présenter la sauvegarde]
Pour l'essentiel de son contenu, elle n'est pas un \emph{filet de sécurité} mais une
\emph{économie}~: elle évite de refaire par le réseau un travail déjà fait. C'est la même
logique que l'exportation, appliquée à la base au lieu des médias. Elle redevient un vrai
filet pour le noyau que le serveur ignore, décrit ci-dessous.
\end{note}

\subsection*{Trois natures de données, trois raisons différentes}

\begin{center}
\begin{tabularx}{\linewidth}{@{}l X l@{}}
\toprule
\thd{Nature} & \thd{Exemples} & \thd{Motif}\\
\midrule
Le serveur l'a aussi &
Messages, conversations, contacts, profil, réactions, listes de contacts, trajets &
Économie\\[4pt]
Seul l'appareil l'a &
Traductions ML~Kit (\code{translatedContent}, \code{sourceLang},
\code{translationState})~; appels de plus de 12~mois~; préférences d'affichage et de
téléchargement &
Filet\\[4pt]
À ne pas restaurer &
Chemins de fichiers locaux, file d'envoi, jetons de session, verrou biométrique &
Nuisible\\
\bottomrule
\end{tabularx}
\end{center}

La deuxième ligne mérite un mot. Les colonnes de traduction portent en commentaire, dans
\code{app\_database.dart}, la mention explicite~: « \emph{ces trois colonnes ne voyagent
jamais~: elles ne sont ni envoyées au serveur ni reçues de lui} ». Un inscrit qui a fait
traduire deux ans de conversations sur son appareil perdrait tout ce travail sans
sauvegarde --- et l'application relancerait l'identification de langue sur tout l'historique
au premier lancement.

L'historique d'appels est le second cas net~: le serveur le purge à 12~mois, l'appareil le
garde. Passé un an, la sauvegarde est le seul endroit où ces appels existent encore.

\subsection*{Les préférences ne sont pas dans la base}

C'est la correction qui justifiait cette section. Les réglages ne vivent pas dans Drift mais
dans \code{SharedPreferences}, un magasin séparé qu'une sauvegarde de la seule base
laisserait entièrement de côté~: langue, thème, échelle de police, réduction des animations,
téléchargement automatique, visibilité des médias, mode d'aperçu des notifications.

Toutes ne doivent pas voyager pour autant. Beaucoup sont déjà côté serveur et reviennent
seules à la connexion~; d'autres sont propres au matériel.

\begin{constat}
\code{PrivacyPrefsService} ne fait que \emph{mettre en cache} la réponse de
\code{getPrivacyPrefs()} sous la clé \code{privacy\_prefs\_json\_v1}. Les préférences de
confidentialité, de notification et les plages « ne pas déranger » ont chacune leur
contrôleur côté serveur. Les sauvegarder les dupliquerait, avec le risque qu'une valeur
périmée écrase la valeur serveur au moment de la restauration.
\end{constat}

Le détail de ce tri figure au chapitre~\ref{ch:sauvegarde-tech}, §\ref{sec:tri}.

\subsection*{Quel poids, réellement}

\begin{alerte}[« Quelques mégaoctets » mérite d'être nuancé]
Le poids de la sauvegarde n'est pas dicté par le nombre de messages mais par les
\textbf{vignettes vidéo}. \code{mediaThumb} stocke une image JPEG encodée en base64 sur la
ligne du message, de l'ordre de 15~Ko pièce. Cinq cents vidéos reçues, et les vignettes
pèsent à elles seules davantage que tout le texte de l'historique.
\end{alerte}

Ordre de grandeur pour un compte chargé --- 20\,000 messages, 500~vidéos --- une fois
compressé~:

\begin{center}
\begin{tabularx}{\linewidth}{@{}X r r@{}}
\toprule
\thd{Poste} & \thd{Brut} & \thd{Après gzip}\\
\midrule
Texte, métadonnées, index & \textasciitilde{}7~Mo & \textasciitilde{}2~Mo\\
Vignettes vidéo (\code{mediaThumb}, base64) & \textasciitilde{}7,5~Mo & \textasciitilde{}5,5~Mo\\
\midrule
\thd{Total} & \textasciitilde{}15~Mo & \textbf{\textasciitilde{}7,5~Mo}\\
\bottomrule
\end{tabularx}
\end{center}

La compression est efficace sur le texte, et rattrape sur les vignettes le tiers de
surcoût que l'encodage base64 avait ajouté. L'estimation reste à confirmer par la mesure~:
\code{StorageInfoService} expose déjà \code{databaseBytes}, il suffit de l'instrumenter sur
quelques comptes réels avant de figer les seuils d'alerte.

\begin{note}[Pourquoi garder les vignettes malgré leur poids]
Elles sont ce qui rend l'historique restauré immédiatement lisible. Les fichiers pleins
arrivent en tâche de fond pendant plusieurs minutes~; sans les vignettes, l'inscrit
contemplerait pendant tout ce temps des colonnes de rectangles gris à la place de ses
souvenirs. Cinq mégaoctets pour que la restauration ait l'air réussie dès la première
seconde, c'est bien acheté.
\end{note}

\section{Où, et sous quelle forme}

Dans le dossier \code{Alanya/} du Drive de l'inscrit --- visible, pas caché. Il voit ses
sauvegardes, connaît leur poids, peut les effacer pour faire de la place. Le champ d'accès
demandé à Google se limite aux fichiers créés par l'application~: Alanya ne voit rien
d'autre de son Drive.

L'archive est \textbf{chiffrée}, avec une clé dérivée du compte Alanya et fournie par le
serveur à la connexion. L'inscrit n'a rien à retenir, et la restauration est transparente.

\begin{alerte}[Ce n'est pas du chiffrement de bout en bout]
La clé étant dérivée côté serveur, Alanya peut techniquement déchiffrer une sauvegarde.
Google, lui, ne le peut pas. C'est un arbitrage assumé contre l'alternative --- un mot de
passe choisi par l'inscrit --- dont le mode de panne (mot de passe oublié = sauvegarde
perdue à jamais) est le premier motif de support sur ce type de fonction. Ce point doit
figurer dans la politique de confidentialité~; le taire serait une faute.
\end{alerte}

\section{Deux sauvegardes, jamais une seule}

On conserve la sauvegarde courante \textbf{et} la précédente. La nouvelle n'écrase l'ancienne
qu'une fois écrite, relue et vérifiée. Le coût est dérisoire --- la base pèse quelques
mégaoctets --- et il ferme le scénario qui rend une sauvegarde inutile~: une écriture
interrompue qui détruit la dernière bonne copie.

\section{Quand elle se déclenche}

La fréquence est au choix de l'inscrit~: \emph{quotidienne}, \emph{hebdomadaire},
\emph{mensuelle}, \emph{jamais}. C'est la convention que WhatsApp a installée et qu'il
saura lire sans explication. Une sauvegarde manuelle reste possible à tout moment, et c'est
elle qu'on met en avant juste avant un changement de téléphone.

\begin{constat}
La condition « seulement en Wi-Fi » est réalisable~: \code{ConnectivityService} distingue
déjà les réseaux non facturés (\code{wifi}, \code{ethernet}) des données mobiles.
\end{constat}

\section{Quand elle échoue}

Quota Google plein, réseau coupé, autorisation retirée~: les échecs passagers sont réessayés
en silence. Ce n'est qu'après \textbf{plusieurs jours consécutifs sans aucune sauvegarde
réussie} que l'application le signale visiblement, avec la cause et la marche à suivre.
Alerter à chaque coupure réseau apprendrait à l'inscrit à ignorer nos alertes --- et il les
ignorerait aussi le jour où elles comptent.

\section{L'inscrit sans compte Google}

Tous n'en ont pas, et certains refuseront d'en connecter un. Ils ne doivent pas rester sans
filet~: l'application écrit alors la même archive dans le stockage du téléphone, sous
\code{Download/Alanya/Sauvegardes/}, à charge pour l'inscrit de la recopier où il veut. Même
format, même chiffrement, même restauration. Seule la destination change.

\section{La restauration sur un nouveau téléphone}
\label{ch:restauration}

À la première connexion sur un appareil neuf, l'application cherche une sauvegarde et
propose de la restaurer. Le déroulement est coupé en deux, parce que les deux moitiés n'ont
pas du tout le même poids.

\begin{enumerate}[leftmargin=*]
  \item \thd{La base, en étape bloquante.} Quelques mégaoctets, quelques secondes, un écran
        dédié. Aucune ambiguïté sur l'état des données~: quand l'inscrit entre dans
        l'application, son historique est là.
  \item \thd{Les médias, en tâche de fond.} Les archives d'export trouvées dans le dossier
        \code{Alanya/} sont rapatriées et redistribuées derrière, reprise après coupure
        comprise, avec un bandeau qui dit où ça en est. L'application est utilisable pendant
        ce temps.
\end{enumerate}

\begin{alerte}[Toujours prévoir la sortie]
Une étape bloquante sans échappatoire enferme l'inscrit dont la sauvegarde est corrompue ou
illisible~: il ne peut plus entrer dans l'application du tout. L'écran de restauration doit
donc toujours porter un « Ignorer et continuer sans restaurer », et l'opération doit rester
relançable ensuite depuis les réglages.
\end{alerte}

\section{Suppression du compte}

Quand l'inscrit supprime son compte Alanya, l'écran de suppression lui \textbf{propose}
d'effacer aussi sa sauvegarde, sans le lui imposer. Ce fichier est dans son espace
personnel~: le détruire d'autorité se défendrait mal, et certains voudront le garder pour
revenir.

% =====================================================================
\part{Comment on le construit}
% =====================================================================

\chapter{Le format d'archive}
\label{ch:format}

\section{Le manifeste}

\begin{lstlisting}[style=Alanya, language={}]
{
  "format": "alanya-export/1",
  "alanyaID": 241030112,
  "generatedAt": "2026-08-30T10:24:07Z",
  "period": { "from": "2026-03-01", "to": "2026-03-31" },
  "filters": { "conversationID": 12, "kind": "all", "owner": "all" },
  "retentionDaysKnown": 30,
  "items": [
    {
      "msgID": 90412, "clientId": "c_8f21ab", "conversationID": 12,
      "conversationName": "Ama Nguema", "senderID": 77, "senderName": "Ama Nguema",
      "isMine": false, "type": 1, "sentAt": "2026-03-14T09:12:44Z",
      "path": "Alanya Images/2026-03-14_ama-nguema_0001.jpg",
      "originalName": "IMG_2231.jpg", "bytes": 2841233,
      "sha256": "9f2c...", "mediaUrl": "/uploads/media/2026-03-14/media_90412_...jpg"
    }
  ],
  "missing": [
    { "msgID": 90455, "sentAt": "2026-03-19T18:02:00Z", "reason": "expired" },
    { "msgID": 90460, "sentAt": "2026-03-21T07:40:11Z", "reason": "neverDownloaded" }
  ]
}
\end{lstlisting}

Trois champs méritent qu'on s'y arrête.

\begin{itemize}[leftmargin=*]
  \item \code{msgID} et \code{clientId} ensemble, jamais l'un sans l'autre. \code{msgID}
        vaut 0 tant que le serveur n'a pas confirmé le message~; \code{clientId} est stable
        dès l'origine. La restauration se cale sur \code{clientId} et n'utilise
        \code{msgID} que comme confirmation.
  \item \code{sha256} rend la restauration \textbf{idempotente}. Rejouer une archive déjà
        restaurée ne duplique rien, et un fichier tronqué par une coupure est détecté au
        lieu d'être installé.
  \item \code{missing} avec son motif~: c'est ce qui transforme une archive incomplète en
        archive honnête.
\end{itemize}

\section{L'assemblage, et pourquoi il ne compresse pas}

\begin{note}[Détail qui change tout sur un export de plusieurs gigaoctets]
Photos JPEG, vidéos MP4, vocaux M4A et PDF sont \textbf{déjà compressés}. Les repasser dans
l'algorithme \emph{deflate} du zip coûte beaucoup de temps processeur et de batterie pour un
gain de taille voisin de zéro --- parfois négatif. L'archive est donc écrite en mode
\emph{stocké} pour ces familles, et compressée uniquement pour le manifeste et l'index HTML.
Le zip ne sert pas à réduire, il sert à \textbf{tenir ensemble}.
\end{note}

L'assemblage lit les fichiers un par un depuis leur emplacement d'origine et les écrit au fil
de l'eau dans l'archive. Aucune copie intermédiaire~: une mise en scène dans un dossier
temporaire doublerait l'espace disque nécessaire, ce qui, sur un téléphone qu'on cherche
précisément à soulager, serait une faute.

\begin{alerte}
Vérifier l'espace libre \emph{avant} de commencer, et refuser proprement si l'archive ne
tient pas. Un export qui remplit le disque à 98~\% laisse l'appareil dans un état pire que
celui d'où l'on partait.
\end{alerte}

\section{Reprise après interruption}

L'export se découpe en trois phases dont la reprise n'a pas le même coût.

\begin{center}
\begin{tabularx}{\linewidth}{@{}l l X@{}}
\toprule
\thd{Phase} & \thd{Ressource} & \thd{Reprise}\\
\midrule
1. Récupération des manquants & réseau &
Gratuite~: les fichiers déjà descendus sont dans le cache, la phase les saute.\\
2. Assemblage du zip & disque et processeur &
Recommencée d'un trait. Sans réseau, c'est rapide~; un zip ne se reprend pas au milieu.\\
3. Dépôt sur Drive & réseau &
Native~: l'API Drive gère les téléversements reprenables.\\
\bottomrule
\end{tabularx}
\end{center}

Les deux phases coûteuses sont donc reprenables, et la seule qui ne l'est pas ne coûte pas
de réseau. C'est ce découpage qui permet de tenir la promesse « tâche de fond reprenable »
sans construire de machinerie de reprise maison.

\section{Le service de premier plan et sa limite de six heures}
\label{sec:fgs}

Un export long tourne dans un service de premier plan. Android encadre désormais ce type de
service, et il faut le regarder en face avant de promettre quoi que ce soit.

\begin{constat}[Vérifié dans la documentation Android]
Les services de premier plan de type \code{dataSync} --- comme \code{mediaProcessing} ---
sont plafonnés à \textbf{six heures par tranche de 24~heures}, pour les applications ciblant
l'API~35 ou plus. Au terme, le système appelle \code{Service.onTimeout()} et laisse quelques
secondes pour appeler \code{stopSelf()}~; passé ce délai il lève
\code{RemoteServiceException}, c'est-à-dire un plantage. Le compteur repart quand l'inscrit
ramène l'application au premier plan. Alanya cible l'API~36~: elle est concernée.
\end{constat}

\begin{note}[Six heures, rapportées à ce qu'on fait, c'est immense]
L'assemblage du zip est local au disque~: des minutes. Le téléversement d'une sauvegarde,
c'est huit mégaoctets. Seul le rapatriement de plusieurs gigaoctets d'archives sur un réseau
lent pourrait en approcher --- et il est déjà reprenable, donc une expiration ne fait que
reporter la suite. Le plafond n'est pas une contrainte pour cette fonctionnalité~; c'en
serait une pour un traitement d'un seul tenant, et c'est une raison de plus de ne pas en
faire un.
\end{note}

Trois exigences suffisent, et aucune n'est coûteuse~:

\begin{enumerate}[leftmargin=*]
  \item Déclarer \code{FOREGROUND\_SERVICE\_DATA\_SYNC} et
        \code{android:foregroundServiceType="dataSync"} dans le manifeste --- qui déclare
        aujourd'hui \code{phone\_call}, \code{microphone}, \code{camera} et \code{location},
        mais pas celui-ci.
  \item Implémenter \code{Service.onTimeout()}~: persister la progression, notifier, appeler
        \code{stopSelf()} \textbf{en quelques secondes}. Ne pas le faire est un plantage, pas
        un avertissement.
  \item Reprendre au tick suivant, ce que le découpage en trois phases permet déjà.
\end{enumerate}

\begin{note}[Et cela se teste, sans attendre six heures]
\begin{lstlisting}[style=Alanya, language={}]
adb shell am compat enable FGS_INTRODUCE_TIME_LIMITS com.alanya237.alanya
adb shell device_config put activity_manager data_sync_fgs_timeout_duration 30000
\end{lstlisting}
Le délai tombe à trente secondes. Le comportement d'expiration devient un test, au lieu d'un
espoir.
\end{note}

\section{Mesurer avant d'annoncer}
\label{sec:mesure}

La feuille d'export annonce un poids avant que l'inscrit ne s'engage. Ce chiffre ne peut pas
venir de la base.

\begin{alerte}[La colonne ment, et pas qu'un peu]
\code{mediaSize} est nullable --- « nul pour les médias envoyés avant que l'app ne relève la
taille des images et vidéos ». Sommer la colonne fait donc compter \textbf{zéro octet} pour
tout média ancien. Le même défaut touchait l'écran \emph{Mes médias}~: son total était
sous-évalué, et surtout son tri « Plus lourds » reléguait les anciens en fin de liste ---
un tri censé faire remonter les gros fichiers cachait précisément les plus anciens.
\end{alerte}

La mesure se fait donc au disque. C'est bon marché~: quelques centaines d'appels \code{stat},
de l'ordre de la milliseconde, et le contrôle d'existence des fichiers en coûtait déjà autant.

\begin{constat}[Correctif déjà appliqué]
\code{ChatDao.setMediaSize(clientId, bytes)} et le rattrapage dans \code{\_onRows} sont
écrits et testés. Un \code{stat} donne l'existence et le poids d'un seul appel~; quand la
colonne est nulle, la taille mesurée est réécrite en base. \textbf{L'écran répare la donnée
en la lisant}~: le total, le tri et les pastilles de poids deviennent justes
définitivement, et l'export héritera d'une base déjà assainie.
\end{constat}

Le même code sert les deux usages --- l'affichage et l'estimation d'export --- pour qu'ils ne
puissent pas diverger.

\chapter{La sauvegarde}
\label{ch:sauvegarde-tech}

\section{Fabriquer l'instantané}

\begin{enumerate}[leftmargin=*]
  \item \code{VACUUM INTO} vers un fichier temporaire. C'est la seule manière d'obtenir un
        instantané \textbf{cohérent} de SQLite sans arrêter l'application~: copier
        \code{talky\_chat.sqlite} pendant qu'une écriture est en cours produirait une base
        corrompue.
  \item \thd{Neutralisation} des colonnes d'état local, \emph{sur la copie} et jamais sur la
        base vivante (§\ref{sec:tri}), et purge des messages fantômes.
  \item Ajout du \code{prefs.json} issu de la liste blanche des préférences.
  \item Compression \emph{gzip} --- ici elle est payante, une base SQLite se comprime bien.
  \item Chiffrement AES-GCM avec la clé du compte.
  \item Téléversement, sous le nom
        \code{alanya-backup-<alanyaID>-<horodatage>.enc}.
  \item Vérification de la taille et de l'empreinte relues, \emph{puis seulement} suppression
        de l'avant-dernière.
\end{enumerate}

Un second fichier, \code{alanya-backup-<alanyaID>-latest.json}, reste \textbf{en clair}~: il
porte la date, la taille, le nombre de messages et la version de schéma. L'écran de réglages
peut ainsi afficher « dernière sauvegarde~: hier, 21~h~04, 8,2~Mo » sans avoir à déchiffrer
quoi que ce soit, ni même à demander la clé au serveur.

\section{Le tri : ce qui monte, ce qui reste}
\label{sec:tri}

\subsection{Les colonnes à neutraliser}

Douze tables composent la base. Aucune n'est exclue en bloc --- ce sont des \emph{colonnes},
au sein de \code{LocalMessages}, qu'il faut vider avant l'instantané. Toutes décrivent l'état
de cet appareil-ci, et les restaurer telles quelles produirait des pannes silencieuses.

\begin{center}
\begin{tabularx}{\linewidth}{@{}l X@{}}
\toprule
\thd{Colonne} & \thd{Pourquoi elle ne doit pas être restaurée}\\
\midrule
\code{localMediaPath} &
Chemin absolu vers \code{media\_cache} sur l'\emph{ancien} téléphone. Restauré tel quel, il
fait croire à l'application que le fichier est là~: la grille \emph{Mes médias} affiche des
éléments qui n'existent pas, et la bulle échoue à l'ouverture.\\[4pt]
\code{pendingUploadPath} &
Fichier en attente d'envoi, resté sur l'ancien appareil. L'outbox tenterait indéfiniment
d'envoyer un fichier introuvable.\\[4pt]
\code{syncPending}, \code{lastEmittedAt}, \code{retryCount}, \code{failureCode} &
État de la file d'envoi. Une sauvegarde prise pendant une coupure réseau ressusciterait, des
mois plus tard, des messages en attente que le serveur a peut-être déjà reçus.\\
\bottomrule
\end{tabularx}
\end{center}

\begin{note}[Le cas \code{localMediaPath} n'est pas qu'un détail d'hygiène]
Il faut le vider \emph{et} laisser la restauration des archives d'export le renseigner à
nouveau, fichier par fichier, à mesure qu'elle réinstalle les médias. C'est le point de
jonction entre les deux mécanismes de ce dossier~: la sauvegarde rétablit les messages, les
archives rétablissent les fichiers, et cette colonne est la couture entre les deux.
\end{note}

Les messages fantômes se traitent au même moment~: \code{ChatDao.purgeGhostMessages} existe
déjà pour les optimistes anciens jamais confirmés. L'appeler avant l'instantané évite de
sauvegarder des lignes que l'application effacerait de toute façon au premier lancement.

\subsection{Les préférences : une liste blanche, jamais une liste noire}

\code{SharedPreferences} contient une quarantaine de clés, dont des jetons d'accès. La
sauvegarde ne doit embarquer que celles qu'on a explicitement désignées.

\begin{alerte}[Ce sens-là et pas l'autre]
Avec une liste noire, toute clé ajoutée plus tard entrerait dans la sauvegarde par défaut ---
y compris un jeton, si quelqu'un en introduit un sans penser à ce document. La liste blanche
fait payer l'oubli du bon côté~: une préférence nouvelle ne sera simplement pas restaurée,
et personne n'y perdra sa sécurité.
\end{alerte}

\begin{center}
\begin{tabularx}{\linewidth}{@{}X X@{}}
\toprule
\thd{Ce qui monte} & \thd{Ce qui reste, et pourquoi}\\
\midrule
\code{app\_locale}, \code{theme\_mode},
\code{app\_font\_scale}, \code{app\_reduce\_motion} &
\code{*\_access\_token}, \code{*\_refresh\_token}, \code{*\_api\_base} --- jetons de session,
sans valeur ailleurs et dangereux à recopier\\[4pt]
\code{auto\_download\_media}, \code{media\_wifi\_only},
\code{media\_data\_saver}, \code{media\_visibility} &
\code{privacy\_prefs\_json\_v1} --- simple cache du serveur, qui le renverra\\[4pt]
\code{notif\_preview\_mode} &
\code{biometric\_lock\_enabled} --- lié au matériel de l'ancien appareil\\[4pt]
\code{call\_ringtone\_selected\_id} (l'identifiant, jamais le chemin) &
\code{call\_ringtone\_active\_path}, \code{call\_ringtone\_native\_name} --- chemins locaux
vers des fichiers absents\\[4pt]
 &
\code{pending\_*}, \code{notif\_dedup\_v1}, \code{ended\_call\_ids\_v1},
\code{status\_seen\_ids}, \code{countries\_cache\_v1} --- files et caches propres à
l'appareil\\[4pt]
 &
\code{media\_retention\_days\_learned} --- réappris au premier \code{410}\\[4pt]
 &
\code{exported\_media\_msg\_ids\_v2} --- journal d'export local, dénué de sens ailleurs\\
\bottomrule
\end{tabularx}
\end{center}

Les préférences retenues voyagent dans un \code{prefs.json} joint à l'instantané, à
l'intérieur de l'archive chiffrée. Les sonneries de liste ne figurent nulle part~: elles sont
déjà synchronisées par compte côté serveur.

\subsection{Le garde-fou : un test qui refuse les clés non classées}

Une liste blanche écrite une fois se périme en silence. Six mois plus tard, quelqu'un ajoute
une préférence, ne connaît pas ce document, et elle n'est simplement jamais restaurée --- sans
que rien ne le signale.

La parade idéale serait un registre \code{PrefKeys} référencé partout. Elle touche les
trente-deux fichiers qui utilisent \code{SharedPreferences}~: trop cher pour ce qu'elle
rapporte ici.

La parade retenue tient en une quarantaine de lignes~: \textbf{un test qui balaie les
sources}. Il collecte les littéraux de clés déclarés dans \code{lib/}, et vérifie que chacun
figure dans exactement un des trois ensembles.

\begin{lstlisting}[style=Alanya, language={}]
lib/core/services/backup/backup_prefs_policy.dart
  kBackedUpKeys     app_locale, theme_mode, app_font_scale, app_reduce_motion,
                    auto_download_media, media_wifi_only, media_data_saver,
                    media_visibility, notif_preview_mode,
                    call_ringtone_selected_id
  kServerOwnedKeys  privacy_prefs_json_v1, ...   -> le serveur les renverra
  kDeviceLocalKeys  *_access_token, biometric_lock_enabled, pending_*, ...
\end{lstlisting}

\begin{note}[Un test qui lit le code source, c'est inhabituel --- et c'est le seul qui tienne]
Un test qui interrogerait \code{SharedPreferences} à l'exécution ne verrait que les clés
déjà écrites sur l'appareil de test, donc jamais celle qu'on vient d'ajouter. Seule
l'inspection des sources voit une clé \emph{déclarée}. Le fichier de classification devient
du même coup la documentation de la règle, et l'échec du test nomme la clé fautive~: la
décision est forcée, pas suggérée.
\end{note}

\section{La clé, et son numéro de version}
\label{sec:cle}

Le serveur expose un point d'accès authentifié qui renvoie une clé dérivée d'un secret
serveur et de l'\code{alanyaID}. Elle n'est jamais écrite en clair sur le téléphone~: on la
demande au moment de sauvegarder, on la demande au moment de restaurer, on l'oublie entre
les deux.

C'est cette dérivation, et non le compte Google, qui rattache une sauvegarde à son
propriétaire~: \textbf{une archive du compte~A posée dans le Drive de~B reste illisible
pour~B}. C'est pourquoi le rattachement peut rester souple.

\subsection{Le jour où le secret doit tourner}

\begin{alerte}[Sans numéro de version, la rotation est une catastrophe]
Fuite, changement d'hébergeur, simple hygiène~: un secret serveur finit toujours par devoir
être remplacé. Si les archives ne portent aucune indication de la clé qui les a chiffrées,
ce remplacement rend \textbf{toutes les sauvegardes existantes illisibles d'un coup}, sans
recours ni avertissement. Posé maintenant, le correctif tient en un champ. Ajouté après
coup, il faut composer avec un parc d'archives anonymes.
\end{alerte}

Chaque archive porte donc un en-tête \textbf{en clair}, devant le contenu chiffré~:

\begin{lstlisting}[style=Alanya, language={}]
ALNB | version de format | kid | algorithme | nonce | alanyaID | <contenu chiffre>
\end{lstlisting}

\code{kid} est le numéro de version de la clé. Il dit \emph{comment} déchiffrer sans qu'il
soit besoin de détenir la clé, ce qui permet de décider avant d'agir.

\subsection{Deux points d'accès, une table qu'on n'efface jamais}

\begin{center}
\begin{tabularx}{\linewidth}{@{}l X@{}}
\toprule
\thd{Point d'accès} & \thd{Rôle}\\
\midrule
\mGET{} \code{/backup/key} & Le \code{kid} courant et sa clé --- pour écrire.\\
\mGET{} \code{/backup/key/:kid} & La clé d'une version donnée --- pour relire une archive
ancienne.\\
\bottomrule
\end{tabularx}
\end{center}

Côté base, une table \code{backup\_key\_secrets(kid, secret, created\_at, retired\_at)} dont
\textbf{aucune ligne n'est jamais supprimée}. Retirer une version signifie uniquement
\emph{cesser d'écrire avec}, jamais \emph{cesser de servir}. La clé effectivement utilisée
est dérivée du secret de la version et de l'\code{alanyaID}, si bien qu'une même version
produit une clé différente pour chaque compte.

La rotation se réduit alors à insérer une ligne et à la marquer courante. Les nouvelles
sauvegardes emploient le nouveau secret, les anciennes restent lisibles. Aucune campagne de
rechiffrement, aucune interruption, aucune archive perdue.

\subsection{Deux compléments à poser en même temps}

\begin{itemize}[leftmargin=*]
  \item \thd{Le \code{kid} figure aussi dans le \code{latest.json} en clair.} L'écran de
        restauration peut ainsi annoncer « sauvegarde illisible, version de clé inconnue »
        \emph{avant} de télécharger huit mégaoctets pour rien.
  \item \thd{Les appels à \code{/backup/key} sont journalisés et limités en débit.} Puisque
        le serveur peut déchiffrer, ce point d'accès est le joyau de la couronne. Une trace
        d'audit fait toute la différence entre « Alanya peut déchiffrer » --- phrase
        inquiétante --- et « Alanya peut déchiffrer, et chaque accès laisse une trace »
        --- phrase défendable devant un inscrit comme devant un régulateur.
\end{itemize}

\section{Rattachement : un Drive par compte Alanya}

La sauvegarde est marquée de l'\code{alanyaID}, pas du compte Google. Se reconnecter depuis
un autre compte Google fonctionne, du moment qu'on ouvre le même compte Alanya~; deux
comptes Alanya sur un même téléphone gardent deux sauvegardes distinctes.

\begin{note}[Pourquoi ne pas figer le couple des deux comptes]
Figer punirait le cas le plus banal~: on change de téléphone, et c'est souvent l'occasion
d'un autre compte Google~; ou bien l'accès à la boîte Gmail est perdu. La sauvegarde
deviendrait irrécupérable, ce qui vide la fonction de son sens. Et le risque que cette
souplesse semble ouvrir --- restaurer la sauvegarde d'autrui --- est déjà fermé par le
chiffrement. Figer n'ajouterait aucune sécurité, seulement un mode de panne.
\end{note}

\section{Restaurer une base plus ancienne que l'application}

La base locale porte une version de schéma (aujourd'hui~26). Deux cas, deux traitements~:

\begin{itemize}[leftmargin=*]
  \item \thd{Sauvegarde plus ancienne que l'application.} Cas normal. Les migrations Drift
        s'exécutent à l'ouverture, comme pour n'importe quelle mise à jour.
  \item \thd{Sauvegarde plus récente que l'application.} L'inscrit a rétrogradé, ou installe
        une version ancienne. La restauration doit \textbf{refuser}, en invitant à mettre
        l'application à jour. Ouvrir une base au schéma futur, c'est la corrompre.
\end{itemize}

\chapter{La restauration}
\label{ch:restauration-tech}

\section{Le vrai risque n'est pas celui qu'on croit}

Une restauration écrit des milliers de lignes dans \code{localMessages}. Or l'application
possède déjà une machinerie qui écrit dans cette même table, à partir du serveur, et qui
part \textbf{toute seule}.

\begin{alerte}[Ce que ferait la conception naïve]
\code{bootstrapEmptyConversationMessages()} se déclenche sans qu'on le lui demande, et sa
condition d'entrée est exactement l'état d'un téléphone neuf. Son commentaire le dit~:
« \emph{précharge les messages des conversations qui ont un aperçu serveur mais aucun message
local confirmé (install fraîche, logout\ldots)} ». Sans précaution, la connexion sur le
nouveau téléphone déclencherait donc le rapatriement de tout l'historique par le réseau ---
pendant que la restauration écrit les mêmes lignes sur la même clé primaire. Il y aurait
course d'écriture~; mais surtout, \textbf{la dépense réseau que la sauvegarde existait pour
éviter aurait lieu quand même}. On aurait payé la sauvegarde pour n'économiser rien.
\end{alerte}

\section{Deux bonnes nouvelles dans le code existant}

\begin{constat}
\textbf{Les curseurs sont déjà justes.} Le bootstrap saute toute conversation qui possède
déjà un message serveur~: \code{if (await \_dao.maxServerMsgId(c.conversID) > 0) continue;}
(\code{conversation\_sync.dart}). Une fois la base restaurée, il ne cible plus rien et se tait
de lui-même. Et \code{syncGlobalDelta} lit \code{\_dao.conversationCursors()}, donc la
synchronisation delta repart naturellement de la crête restaurée. \textbf{Aucune plomberie de
curseur n'est à construire.}
\end{constat}

\begin{constat}
\textbf{La séquence de démarrage est linéaire et le socket arrive en dernier.} Dans
\code{main.dart}, l'enchaînement après authentification est ordonné, et le socket n'est
connecté qu'après les \code{bind} --- un commentaire du fichier explique pourquoi. Chaque
chemin de synchronisation est par ailleurs gardé par \code{if (\_myId == 0) return;}, et
\code{\_myId} n'est posé que par \code{ChatRepository.bind()}.
\end{constat}

\section{Quand demande-t-on le compte Google ?}
\label{sec:quand-google}

Proposer une restauration suppose de savoir qu'une sauvegarde existe. Chercher dans Drive
imposerait de connecter Google \emph{avant} que l'inscrit soit entré dans l'application ---
au tout premier démarrage, sur un téléphone qu'il vient d'allumer. C'est le pire moment pour
réclamer un compte tiers, et on le réclamerait aussi à ceux qui n'ont jamais rien
sauvegardé.

\textbf{La réponse est de ne pas interroger Drive du tout à ce stade}~: le serveur peut le
savoir à sa place.

À chaque sauvegarde réussie, l'application déclare au serveur de simples métadonnées --- date,
taille, \code{kid}, et l'adresse Google \textbf{masquée}. Aucun contenu, jamais. À la
connexion sur un appareil neuf, le serveur les rend, et l'application sait qu'une sauvegarde
existe sans avoir touché à Google.

\begin{lstlisting}[style=Alanya, language={}]
identifiants verifies
  -> le serveur annonce : sauvegarde 29/08, 8,2 Mo
     compte a...@gmail.com
  -> "Restaurer" / "Ignorer"      <- Google demande ici
  -> connexion Google + lecture Drive
  -> restauration (bloquante)
  -> chatProvider.bind(myId)
  -> statusProvider.bind(myId)
  -> connexion du socket
\end{lstlisting}

Trois gains, et le troisième est le moins évident~:

\begin{itemize}[leftmargin=*]
  \item Un inscrit véritablement nouveau \textbf{n'est jamais sollicité} pour un compte
        Google.
  \item L'écran ne pose plus une question spéculative~; il annonce un fait daté et chiffré,
        ce qui rend le choix « Restaurer / Ignorer » réellement éclairé.
  \item Le cas « mauvais compte Google » cesse d'être une impasse. Sans cette métadonnée,
        l'application ne sait dire que « aucune sauvegarde trouvée » --- indiscernable d'une
        absence réelle. Avec elle~: « une sauvegarde existe, connectez-vous avec
        \code{a\ldots@gmail.com} ».
\end{itemize}

\begin{alerte}[L'adresse est stockée masquée, jamais en clair]
C'est une donnée personnelle dont nous n'avons aucun besoin complet~: la première lettre et
le domaine suffisent à guider l'inscrit vers le bon compte. En conserver davantage serait
accumuler du risque pour rien.
\end{alerte}

Il faut enfin garder une entrée manuelle « restaurer une sauvegarde » dans les réglages~:
pour les sauvegardes antérieures à cette métadonnée, et pour le cas où le serveur l'aurait
perdue.

\section{La solution : un ordre, pas un verrou}

\begin{lstlisting}[style=Alanya, language={}]
identifiants verifies
  -> ecran de restauration (bloquant)      <- le point d'insertion
  -> chatProvider.bind(myId)
  -> statusProvider.bind(myId)
  -> connexion du socket
\end{lstlisting}

Rien ne court encore, donc rien ne peut se percuter. Il n'y a ni verrou à poser, ni file
d'attente à gérer, ni état concurrent à arbitrer.

\begin{note}[Une décision produit qui paie une dette technique]
C'est le choix d'une \emph{restauration bloquante} qui rend cette solution possible. Une
restauration en arrière-plan aurait imposé un verrou, une reprise de curseur et un arbitrage
entre deux écrivains. La contrepartie --- quelques secondes d'attente au premier démarrage
--- achète la disparition complète d'une classe de bugs.
\end{note}

\section{L'état de restauration, et pourquoi il est indispensable}

C'est la seule pièce réellement nouvelle~: un état persisté hors de la base, puisque la base
est justement ce qu'on est en train de remplacer.

\begin{center}
\begin{tabularx}{\linewidth}{@{}l X@{}}
\toprule
\thd{État} & \thd{Signification}\\
\midrule
\code{inconnu} & On n'a pas encore cherché de sauvegarde.\\
\code{proposée} & Une sauvegarde existe, l'inscrit n'a pas tranché.\\
\code{en cours} & Écriture commencée. \textbf{La base est dans un état intermédiaire.}\\
\code{faite} & Terminée et vérifiée.\\
\code{ignorée} & L'inscrit a refusé. On ne repropose plus au démarrage.\\
\bottomrule
\end{tabularx}
\end{center}

\begin{alerte}[Une demi-restauration est pire que pas de restauration]
Si l'application est tuée au milieu de l'écriture, la base contient des messages pour
\emph{certaines} conversations seulement. Au redémarrage, le bootstrap les considère comme
complètes --- puisque \code{maxServerMsgId > 0} --- et les saute définitivement.
\textbf{Les trous deviennent permanents, et silencieux.} C'est précisément parce que la
logique de curseur est correcte qu'un état intermédiaire est toxique. La règle est donc
sans nuance~: trouver l'état \code{en cours} au démarrage impose de vider la base et de
reproposer la restauration depuis le début.
\end{alerte}

\section{Les cinq répétitions à écrire avant le code}

Une restauration ne se vérifie pas en la relisant. Ces cinq scénarios doivent être des tests,
et ils doivent exister avant l'implémentation qu'ils décrivent.

\begin{center}
\begin{tabularx}{\linewidth}{@{}l X@{}}
\toprule
\thd{Scénario} & \thd{Comportement attendu}\\
\midrule
Archive tronquée & Détectée à la vérification d'intégrité. La base n'est pas touchée,
l'inscrit peut réessayer ou ignorer.\\
Schéma plus récent que l'app & Refus explicite, invitation à mettre à jour. Jamais
d'ouverture d'une base au schéma futur.\\
\code{kid} inconnu & Message clair avant tout téléchargement, grâce au \code{latest.json}
(§\ref{sec:cle}).\\
Coupure en pleine écriture & Au redémarrage, état \code{en cours} détecté~: purge et
nouvelle proposition.\\
Aucune sauvegarde trouvée & Distinguer « ce compte Google n'en contient pas » de « erreur
réseau ». Proposer d'essayer un autre compte Google dans le premier cas.\\
\bottomrule
\end{tabularx}
\end{center}

\begin{note}[Toujours une porte de sortie]
Chacun de ces écrans doit porter « Ignorer et continuer sans restaurer », et l'opération doit
rester relançable depuis les réglages. Une étape bloquante sans échappatoire enferme dehors
l'inscrit dont la sauvegarde est illisible.
\end{note}

\subsection{Ces répétitions imposent une contrainte d'architecture}

Écrire ces cinq scénarios est facile~; les rendre \emph{jouables} ne l'est que si on le
décide avant de coder.

\begin{alerte}[La règle qui rend le plan de test possible]
\code{RestoreService} reçoit son \code{BackupTarget} et son fichier source \textbf{en
paramètres}. Il ne va jamais les chercher dans un singleton, ni ouvrir Drive lui-même. Sans
cette règle, les cinq scénarios ne seraient rejouables qu'à la main sur un vrai téléphone
avec un vrai compte Google --- autant dire jamais.
\end{alerte}

Avec elle, un \code{FakeBackupTarget} sert des archives fabriquées, et chaque scénario tient
en quelques lignes~:

\begin{center}
\begin{tabularx}{\linewidth}{@{}l X@{}}
\toprule
\thd{Fixture} & \thd{Fabrication}\\
\midrule
Archive tronquée & Une archive valide, coupée aux trois quarts.\\
\code{kid} inconnu & En-tête valide, \code{kid} inexistant côté serveur.\\
Schéma futur & Base au \code{schemaVersion} 99.\\
Mauvais compte & Archive d'un autre \code{alanyaID}, donc indéchiffrable.\\
Cible vide & \code{FakeBackupTarget} qui renvoie une liste vide.\\
Coupure & Interruption injectée au milieu de l'écriture~; on vérifie qu'un
redémarrage trouve l'état \code{en cours} et purge.\\
\bottomrule
\end{tabularx}
\end{center}

\begin{constat}
Le socle de test existe déjà. \code{AppDatabase.forTesting(NativeDatabase.memory())} est
utilisé par la suite actuelle~: restaurer dans une base en mémoire et vérifier son contenu
est directement à portée, sans émulateur ni appareil.
\end{constat}

\chapter{Ce qui reste à lever}
\label{ch:a-lever}

\section{L'inconnue qui conditionne l'architecture}

Toute la restauration repose sur une hypothèse non vérifiée~: sur un téléphone neuf,
l'application retrouve dans le Drive de l'inscrit les fichiers qu'elle y avait déposés depuis
l'ancien. Le champ d'accès retenu limite l'application aux fichiers qu'elle a créés --- reste
à savoir si ce droit est rattaché à \textbf{l'application} (identifiant client OAuth) ou à
\textbf{l'installation}.

\begin{center}
\begin{tabularx}{\linewidth}{@{}l X@{}}
\toprule
\thd{Si le droit suit\ldots} & \thd{Conséquence}\\
\midrule
l'application & La sauvegarde est visible sur le nouveau téléphone. La conception fonctionne
telle qu'écrite.\\
l'installation & L'application est aveugle à une archive pourtant présente. La restauration
automatique disparaît.\\
\bottomrule
\end{tabularx}
\end{center}

Les trois issues en cas de réponse défavorable, et aucune n'est indolore~: basculer sur le
dossier caché (on perd le dossier visible que l'inscrit peut gérer, et la vérification Google
s'alourdit), faire choisir le fichier à la main par le sélecteur Google (la restauration
automatique disparaît du parcours, au pire moment --- le premier démarrage d'un téléphone
neuf), ou demander l'accès complet au Drive (audit annuel payant, hors de question).

\subsection{Isoler d'abord}

Toute la conception ne touche Drive que par quatre gestes~: lister, écrire, lire, supprimer.
Derrière une interface \code{BackupTarget}, la réponse au test ne coûte plus qu'une classe~:

\begin{lstlisting}[style=Alanya, language={}]
abstract class BackupTarget {
  Future<List<RemoteArchive>> list();
  Future<void> write(File local, String name);
  Future<File> read(String id, File into);
  Future<void> delete(String id);
}
    DriveVisibleFolderTarget   <- choix par defaut
    DriveAppDataTarget         <- repli si le test echoue
    LocalFolderTarget          <- deja requis (sans compte Google)
\end{lstlisting}

\begin{note}[Cette abstraction ne coûte rien de plus]
La décision prise sur les inscrits sans compte Google imposait déjà \code{LocalFolderTarget}.
L'interface devait donc exister de toute façon~: elle protège du risque Drive par pure
retombée.
\end{note}

\subsection{Tester ensuite --- et deux variantes, jamais une seule}

\begin{enumerate}[leftmargin=*]
  \item Déclarer l'identifiant client OAuth dans la console Google.
  \item Depuis un simple script --- Flutter n'est pas nécessaire --- s'authentifier avec un
        compte de test et créer \code{Alanya/test.txt}.
  \item Rejouer l'accès de deux façons~:
    \begin{itemize}
      \item \thd{Réinstallation simple}~: repartir d'un jeton vierge sans rien révoquer.
            C'est le cas réel le plus fréquent.
      \item \thd{Révocation puis ré-autorisation}~: l'inscrit a retiré l'accès depuis son
            compte Google, ou trop de temps a passé.
    \end{itemize}
  \item Redemander la liste dans chacun des deux cas.
\end{enumerate}

\begin{alerte}[Ne pas se contenter de la première variante]
Les deux peuvent parfaitement donner des réponses différentes~: une révocation est
susceptible de jeter les droits accordés fichier par fichier, là où une réinstallation les
conserverait. N'éprouver que la réinstallation donnerait une fausse assurance.
\end{alerte}

Coût~: une demi-journée. À faire \textbf{avant} la première ligne du reste.

\section{Les points désormais tranchés}

Cette section listait cinq questions ouvertes. Elles ont toutes reçu une réponse, consignée
au fil du document~; on n'en garde ici que la table de renvoi, pour que rien ne se perde.

\begin{center}
\begin{tabularx}{\linewidth}{@{}l X l@{}}
\toprule
\thd{Sujet} & \thd{Réponse} & \thd{Où}\\
\midrule
Service de premier plan &
Type \code{dataSync} déclaré, \code{onTimeout()} implémenté, travail reprenable. Le plafond
de six heures n'est pas une contrainte pour ce que nous faisons. &
§\ref{sec:fgs}\\[4pt]
Poids annoncé avant export &
Mesure au disque et rattrapage en base. \textbf{Correctif déjà appliqué} au code de
« Mes médias ». &
§\ref{sec:mesure}\\[4pt]
Moment de la connexion Google &
Le serveur porte la métadonnée de sauvegarde~: Google n'est demandé qu'après un « Restaurer »
explicite. &
§\ref{sec:quand-google}\\[4pt]
Plan de test &
Six fixtures, et la règle d'injection qui les rend jouables. &
§\ref{ch:restauration-tech}\\[4pt]
Garde-fou de la liste blanche &
Un test qui balaie les sources et refuse toute clé non classée. &
§\ref{sec:tri}\\
\bottomrule
\end{tabularx}
\end{center}

\begin{note}[Le seul point encore ouvert de ce chapitre]
Reste la vérification que l'exportation hérite bien de l'exclusion des médias à vue unique.
\code{watchLocalMedia} les écarte déjà côté destinataire et les conserve côté expéditeur~;
il faut confirmer que ce comportement est celui qu'on veut dans une archive en clair.
C'est une relecture de dix minutes, pas une décision d'architecture.
\end{note}

\chapter{Impact sur le code}
\label{ch:impact}

\section{Ce qui existe et se réutilise}

\begin{center}
\begin{tabularx}{\linewidth}{@{}l X@{}}
\toprule
\thd{Existant} & \thd{Rôle dans cette conception}\\
\midrule
\code{ChatDao.watchLocalMedia} & Fournit déjà le périmètre exact à exporter, filtres compris.\\
\code{AlanyaMediaExportService} & Arborescence des dossiers, album galerie, journal des
identifiants déjà exportés.\\
\code{MediaExpiryPolicy} & Distingue « récupérable » de « purgé » sans tenter la requête.\\
\code{ConnectivityService} & Condition Wi-Fi de la sauvegarde (\code{wifi}, \code{ethernet}).\\
\code{StorageInfoService} & Espace libre avant export, poids de la base pour la sauvegarde.\\
\code{flutter\_local\_notifications} & Notification de progression.\\
\code{CallMediaForegroundService} & Motif de service de premier plan déjà maîtrisé, à
reprendre pour l'export long.\\
\bottomrule
\end{tabularx}
\end{center}

\section{Ce qu'il faut ajouter}

\begin{center}
\begin{tabularx}{\linewidth}{@{}l X@{}}
\toprule
\thd{Dépendance} & \thd{Pourquoi}\\
\midrule
\code{google\_sign\_in} & Connexion au compte Google.\\
\code{googleapis} (Drive~v3) & Dépôt, listage, téléversement reprenable.\\
\code{archive} & Écriture du zip au fil de l'eau, mode stocké.\\
\code{cryptography} & AES-GCM.\\
\code{workmanager} & Planification périodique de la sauvegarde.\\
\bottomrule
\end{tabularx}
\end{center}

\begin{alerte}[Démarche Google à ne pas découvrir la veille de la publication]
Le champ d'accès retenu se limite aux fichiers créés par l'application, c'est le moins
exigeant de ceux qui conviennent --- nettement plus léger que l'accès complet au Drive, qui
impose un audit de sécurité annuel et payant. Il n'en reste pas moins qu'un écran de
consentement OAuth doit être publié et que l'application doit être vérifiée par Google avant
mise en production. À vérifier contre la politique Google en vigueur au moment du
développement, et à lancer \textbf{tôt}~: le délai n'est pas maîtrisé par l'équipe.
\end{alerte}

\section{Fichiers à créer}

\begin{lstlisting}[style=Alanya, language={}]
lib/core/services/export/
  media_export_job.dart        phases, progression, reprise
  export_manifest.dart         ecriture et relecture du manifeste
  export_zip_builder.dart      assemblage au fil de l'eau
  export_index_html.dart       planche-contact

lib/core/services/backup/
  backup_service.dart          orchestration
  backup_crypto.dart           AES-GCM, en-tete ALNB, kid
  backup_scheduler.dart        frequence, conditions, reessais
  backup_snapshot.dart         VACUUM INTO, neutralisation, prefs.json
  backup_target.dart           interface : list / write / read / delete
  drive_visible_folder_target.dart   choix par defaut
  drive_appdata_target.dart          repli si le test Drive echoue
  local_folder_target.dart           inscrit sans compte Google
  restore_service.dart         base bloquante, medias en fond
  restore_state.dart           inconnu / proposee / en cours / faite / ignoree

lib/screens/profile/
  backup_screen.dart           reglages et etat
  restore_screen.dart          etape bloquante, avec sortie
  my_media_screen.dart         + bouton et feuille d'export

lib/main.dart                  insertion de l'etape de restauration
                               AVANT chatProvider.bind(myId)
\end{lstlisting}

\section{Côté serveur}

Deux points d'accès et une table, tous décrits au §\ref{sec:cle}~:
\mGET{} \code{/backup/key} pour écrire, \mGET{} \code{/backup/key/:kid} pour relire une
archive ancienne, et \code{backup\_key\_secrets} dont aucune ligne n'est jamais supprimée.
Les deux points d'accès sont authentifiés comme le reste de l'API, journalisés et limités en
débit.

La rétention, elle, n'appelle aucun développement~: elle est déjà annoncée par le champ
\code{retentionDays} des réponses \code{410}.

\chapter{Décisions arrêtées}
\label{ch:decisions}

\begin{center}
\begin{tabularx}{\linewidth}{@{}l X@{}}
\toprule
\thd{Sujet} & \thd{Décision}\\
\midrule
Bouton Exporter & Archive zip assemblée localement, remise au partage système, à
Téléchargements, ou déposée sur Drive.\\
Contenu exporté & Médias de la période, plus \code{manifest.json} et \code{index.html}.
Pas les messages texte --- v2 possible.\\
Éléments manquants & Récupérables et purgés distingués~; téléchargement proposé, chiffré
d'avance, non coché par défaut.\\
Contenu sauvegardé & La base, moins les colonnes d'état local, plus une liste blanche de
préférences. Médias exclus.\\
Tri des préférences & Liste blanche explicite. Jamais une liste noire.\\
Clé de sauvegarde & Numérotée (\code{kid}) dans un en-tête en clair. Table de secrets dont
aucune ligne n'est supprimée. Accès journalisé.\\
Ordre au démarrage & La restauration s'insère \emph{avant} \code{chatProvider.bind(myId)}.
Un ordre, pas un verrou.\\
État de restauration & Persisté hors base. \code{en cours} au démarrage $\Rightarrow$ purge
et nouvelle proposition.\\
Accès au stockage & Interface \code{BackupTarget}, trois implémentations. Le résultat du
test Drive ne coûte qu'une classe.\\
Emplacement & Dossier \code{Alanya} visible du Drive de l'inscrit.\\
Chiffrement sauvegarde & AES-GCM, clé dérivée du compte côté serveur. Pas de bout en bout,
et on le dit.\\
Chiffrement archive export & Aucun. Elle doit s'ouvrir sans Alanya. L'écran avertit.\\
Rattachement & Un Drive par compte Alanya. Compte Google non figé.\\
Fréquence & Au choix~: quotidienne, hebdomadaire, mensuelle, jamais. Manuelle toujours
possible.\\
Versions conservées & Deux~: la courante et la précédente. Écrasement seulement après
vérification.\\
Échecs & Réessai silencieux, alerte visible après plusieurs jours sans succès.\\
Sans compte Google & Même archive écrite dans \code{Download/Alanya/Sauvegardes/}.\\
Restauration & Base en étape bloquante avec sortie possible~; médias en tâche de fond
reprenable.\\
Suppression de compte & Effacement de la sauvegarde proposé, jamais imposé.\\
Badge de rétention & Valeur apprise du serveur, formulée par contraste~: sans limite sur
l'appareil, 30~jours sur les serveurs.\\
Plateforme & Android d'abord. Abstraction de destination prête pour iCloud.\\
\bottomrule
\end{tabularx}
\end{center}

\section{Hors périmètre de ce lot}

\begin{itemize}[leftmargin=*]
  \item Le menu \textbf{Plus} et la carte \textbf{« Il y a un an, ce jour-là »} de la
        maquette~: fonction d'engagement, conception séparée.
  \item \textbf{iOS et iCloud}~: l'abstraction de destination est prévue, l'implémentation
        ne l'est pas.
  \item Le \textbf{transfert direct d'appareil à appareil} en Wi-Fi local. C'est pourtant la
        réponse la plus fidèle à l'esprit de la note --- aucun octet ni chez nous ni chez
        Google. À reconsidérer si le dépôt d'archives sur Drive se révèle trop peu utilisé.
  \item L'export du \textbf{fil de discussion lisible} en HTML.
\end{itemize}

\section{Dépendances : état réel}

Cette section listait quatre préalables. Vérification faite dans le code, deux étaient déjà
levés.

\begin{constat}
\code{evictIfNeeded()} et le plafond de 500~Mo ont été \textbf{retirés} de
\code{media\_cache\_service.dart}. Le commentaire du fichier porte le raisonnement~: évincé
côté appareil \emph{et} tombé côté serveur = perdu définitivement~; les deux mécanismes
étaient anodins séparément et destructeurs ensemble. La promesse « sans limite sur
l'appareil » est donc tenue.
\end{constat}

\begin{constat}
Le \code{410 MEDIA\_EXPIRED} est \textbf{géré}. \code{MediaExpiryPolicy.noteResponse()}
intercepte la réponse dans le chemin de téléchargement et mémorise la rétention annoncée~;
\code{isExpired()} évite même d'émettre la requête. La distinction « récupérable » /
« purgé » de la feuille d'export est donc réalisable dès maintenant.
\end{constat}

Il reste deux points, et un seul concerne le code~:

\begin{enumerate}[leftmargin=*]
  \item \thd{Le test de visibilité Drive} (§\ref{ch:a-lever}, deux variantes, une
        demi-journée). Il ne bloque que le choix entre \code{DriveVisibleFolderTarget} et
        \code{DriveAppDataTarget} --- une classe derrière l'interface \code{BackupTarget}.
        Tout le reste peut être écrit sans attendre sa réponse.
  \item \thd{L'écran de consentement OAuth publié et l'application vérifiée par Google.}
        Ce n'est pas un préalable au code mais au \emph{déploiement}, et le délai n'est pas
        maîtrisé par l'équipe~: à lancer le premier jour, en parallèle du développement.
\end{enumerate}

\begin{note}[Ce qui n'est plus une inconnue]
La course entre restauration et synchronisation s'est révélée être un simple problème
d'ordre~: les curseurs existants sont déjà justes et la séquence de démarrage est linéaire
(chapitre~\ref{ch:restauration-tech}).
\end{note}

\vfill
\begin{center}\small\color{AlanyaGray}\sffamily
Fin du dossier --- exportation, sauvegarde et restauration.\\
À valider avant écriture de code.
\end{center}

\end{document}
